\section*{第 \S\arabic{exercisegroup} 节习题}
\phantomsection
\addcontentsline{toc}{section}{第 \protect\S\arabic{exercisegroup} 节习题}

\begin{enumerate}
\def\labelenumi{\arabic{enumi})}
\tightlist
\item
  证明：对实变量 \(x\) （趋于 \(+\infty\) ），函数
\end{enumerate}

\[
f (x) = \left(x \cos^ {2} x + \sin^ {2} x\right) e ^ {x ^ {2}}
\]

单调，但它既不与 \(e^{x^2}\) 可比较，也不与 \(x^{1/2}e^{x^2}\) 弱可比较。

\begin{enumerate}
\def\labelenumi{\arabic{enumi})}
\setcounter{enumi}{1}
\tightlist
\item
  设 \(\varphi\) 是一个严格正的函数，对 \(x > 0\) 有定义且递增，并满足 \(\varphi \gg 1\) 。
\end{enumerate}

\begin{enumerate}
\def\labelenumi{\alph{enumi})}
\item
  证明：若函数 \(\log \varphi(x) / \log x\) 递增，则函数之间的关系 \(f \ll g\) （这些函数 \(>0\) ）蕴含 \(\varphi \circ f \ll \varphi \circ g\) （当 \(g \gg 1\) 时）。
\item
  给出一个例子，其中 \(\log \varphi(x)/\log x\) 递减，f 与 g 是两个 \textgreater0 的函数，满足 \(g \gg 1\) 且 \(f \sim g\) ，但 \(\varphi \circ f\) 不等价于 \(\varphi \circ g\) 。
\end{enumerate}

3) a) 设 \((\alpha_{i},\beta_{i})\) 是 \(n\) 个互异的实数对（由实数 \(\geqslant 0\) 组成），异于 \((0,0)\) ，且满足 \(\inf_{i}\alpha_{i} = \inf_{i}\beta_{i} = 0\) 。考察方程

\[
f (x, y) = \sum_ {i = 1} ^ {n} a _ {i} x ^ {\alpha_ {i}} y ^ {\beta_ {i}} \left(1 + \varphi_ {i} (x, y)\right) = 0
\]

其中诸 \(a_{i}\) 是 \(\neq0\) 的实数，诸 \(\varphi_{i}\) 是正方形 \(0\leqslant x\leqslant a,0\leqslant y\leqslant a\) 上的连续函数，且当 \((x,y)\) 趋于 \((0,0)\) 时趋于 0。假设存在一个函数 g，它在区间 {[}0,b{]} 上为正且连续，随 x 趋于 0，且满足 \(f(x,g(x))=0\) （对所有 \(x\in[0,b]\) ）。证明：对每个实数 \(\mu>0\) ，当 x 趋于 0 时 \(g(x)/x^{\mu}\) 趋于一个有限或无穷的极限（利用 V, p.~217, prop. 9，对每个数 \(t\geqslant0\) 考察方程 \(f(x,tx^{\mu})=0\) ）。

\begin{enumerate}
\def\labelenumi{\alph{enumi})}
\setcounter{enumi}{1}
\item
  为使 \(g(x)/x^{\mu}\) 趋于一个有限极限 \(\neq0\) ，必须 \(\mu\) 满足：至少存在两个互异的对 \((\alpha_{h},\beta_{h})\) 与 \((\alpha_{k},\beta_{k})\) ，使 \(\alpha_{h}+\mu\beta_{h}=\alpha_{k}+\mu\beta_{k}\) ，且对每个其余的对 \((\alpha_{i},\beta_{i})\) 有 \(\alpha_{i}+\mu\beta_{i}\geqslant\alpha_{h}+\mu\beta_{h}\) 。这样就得到有限多个可能的值 \(\mu_{j}(1\leqslant j\leqslant r)\) ；数 \(-1/\mu_{j}\) 是平面 \(R^{2}\) 中这样的仿射直线的斜率：这些直线至少包含点集 \((\alpha_{i},\beta_{i})\) 中的两个点，且该点集的其余各点都位于所考虑的直线上方（“牛顿多边形”）。
\item
  设 \(\mu_{1}\) 是诸数 \(\mu_{j}\) 中最小者。证明 \(g(x) / x^{\mu_1}\) 趋于一个有限极限（可能为零）。（首先证明总可假定：若 \(i\) 与 \(j\) 是两个互异的指标，则不能同时有 \(\alpha_{i} \leqslant \alpha_{j}\) 与 \(\beta_{i} \leqslant \beta_{j}\) ；由此推出可设 \(\alpha_{1} = 0\) ， \(\alpha_{i} > 0\) 与 \(\beta_{i} < \beta_{1}\) （对 \(i \neq 1\) ）；然后令 \(g(x) = t(x)x^{\mu_1}\) ，证明 \(t(x)\) 不能趋于 \(+\infty\) （当 \(x\) 趋于 0）。）
\end{enumerate}

\(d)\) 由 \(c\) ) 出发，对 \(r\) 作递归，推出存在一个指标 \(j\) ，使 \(g(x) / x^{\mu_j}\) 趋于一个有限极限 \(\neq 0\) （当 \(x\) 趋于 0）。

\setcounter{exercisegroup}{2}
\refstepcounter{exercisegroup}
